% Discussion

\subsection{Model-based vs binary mask based diameter measurement}
\label{subsec:model_vs_binary}
The model-based and binary mask based diameter methods define vessel boundary in fundamentally different ways. Binarisation of the en face maps discards the spatial intensity distribution for a simple binarisation and reduces each pixel to a binary decision. The measured diameter directly depends on the chosen threshold value and on angiogram intensity, both of which may vary across imaging sessions, contrast conditions, and even between vessels within the same scan. Model-based methods fit a parametric model function to the maximum or mean intensity profiles, and calculate the diameter from the fitted parameters. As the fit captures the shape of the profile rather than the absolute intensity, it is inherently less sensitive to changes in angiogram signal and does not require a threshold at all.

A practical consequence of this difference is how correctable each method is. Because the model-based fit requires no threshold, it returns a single consistent value for a given profile, whereas the binarisation slope shifts with the chosen threshold, so any calibration must also encode which threshold was used. Underestimation can be corrected by a fixed proportion using the slope, provided the relationship with the ground truth is strongly linear, but such a rescaling removes only the systematic offset and leaves the scatter about the regression unchanged. Linearity, rather than slope alone, therefore determines how reliably a method can be corrected. Post-contrast every model-based metric reached $R^2 \geq 0.92$, whereas binary mask based MIP fell from 0.81 at BW-1 to 0.48 at BW-5, so more than half the variance would remain after any correction. This is what makes the higher-linearity model-based metrics usable even pre-contrast, where $w_{1/e^2,\mathrm{GG}}$ retained $R^2 = 0.81$ against 0.62 for the best binary mask based method.

This difference explains the trends observed in Fig.~\ref{fig:cov_r2} and Table~\ref{tab:accuracy_linearity}. Post-contrast, every model-based metric was more linear than every binary mask based method, and the GG $1/e^2$ metric reached near-unity slopes at peak and post-contrast without a threshold parameter to tune, whereas for the binary mask based methods linearity was overall lower and depended on threshold. Because the peak angiogram intensity varies between vessels, a fixed threshold removes a different fraction of the cross-section from each one. The measured diameters therefore vary from the ground truth by different amounts, which weakens the linearity. Pre-contrast, the hourglass artefact caused by RBC orientation within vessel lumen suppresses edge signals, causing all methods to underestimate diameter~\cite{bernucciInvestigationArtifactsRetinal2018}. However, model-based methods still retained more of that linearity, producing systematic, predictable underestimation. Binary mask based methods on the other hand showed weaker linearity in most conditions, potentially because the suppressed edge signals did not consistently exceed the binarisation threshold, and because different vessels have different intensity.

The same mechanism accounts for how each method responded to the change in angiogram intensity from peak to post-contrast. Angiogram intensity is highest as the bolus arrives and then settles to a lower steady state (Fig.~\ref{fig:dycoct}f), so a fixed threshold encloses a smaller fraction of each cross-section post-contrast than at peak, an effect visible in the binarised cross-sections over time (Fig.~\ref{fig:dycoct}d,e). Model-based metrics were largely insensitive to this, changing by at most 0.04 in slope and 0.04 in $R^2$ between the two conditions. Binary mask based measurements from MIP degraded progressively with threshold, the slope falling by 0.04 at BW-1 but 0.11 at BW-5, and $R^2$ by 0.07 and 0.30 respectively, while MeIP remained comparatively stable (slope change 0.02--0.03 at every threshold). Model-based methods therefore showed comparable agreement with the ground truth at both peak and post-contrast, whereas the agreement of binarisation depended on the combination of threshold and contrast condition.

The generalised Gaussian (GG) model will always outperform the standard Gaussian because the intensity profiles of the vessels are usually not purely Gaussian (Fig.~\ref{fig:gg_model}b). The standard Gaussian cannot capture relatively flat-topped angiogram intensity profiles and tends to fit a narrower curve, whereas the GG model accommodates through the shape parameter $\beta$, which ranges from Gaussian ($\beta = 2$) to increasingly flat-topped profiles ($\beta > 2$). The flat-topped shapes arise because the measured profile is effectively the vessel's reflectivity profile convolved with the system point spread function. Because the GG reduces to a standard Gaussian when $\beta = 2$ and departs from it only when the data support a flatter profile, a correctly applied GG fit should perform at least as well as the Gaussian, at the cost of additional computation.

Among the width metrics, $w_{fwhm,\,\mathrm{GG}}$ consistently showed the weakest accuracy, although its linearity remained high post-contrast ($R^2$ = 0.96 for MIP and 0.98 for MeIP), indicating that the underestimation is systematic and therefore correctable. Near the half-maximum the profile shape changes rapidly and varies between vessels, from Gaussian to flat-topped, so the width measured at that level is more variable. The $1/e^2$ width captures more of the transition region and so may have provided a more accurate measurement of the vessel diameter.

\subsection{Maximum vs mean intensity projection}
\label{subsec:max_vs_mean}
Across nearly all methods and conditions, MeIP outperformed MIP in both accuracy and precision (Table~\ref{tab:accuracy_linearity}). Linearity was the exception. Pre-contrast MIP produced higher $R^2$ than MeIP for all model-based metrics, and MIP also produced higher $R^2$ at the two lowest binarisation thresholds, with MeIP outperforming MIP only post-contrast.  This finding is noteworthy because MIP has been shown to produce higher signal-to-noise ratio (SNR) and contrast in en face OCTA images~\cite{hormelMaximumValueProjection2018}, and is the default in many commercial OCTA systems and analysis pipelines~\cite{sampsonTowardsStandardizingRetinal2022,untrachtTowardsStandardisingRetinal2024}. However, we observed that higher SNR did not necessarily translate to more accurate diameter measurements. The MIP selects the single highest intensity value along each A-scan within the vessel ROI, which makes it inherently sensitive to outliers. MeIP averages all intensity values along each A-scan, which smooths out intensity fluctuations and provides a consistent profile. This averaging could be beneficial for vessel diameter measurement, as the goal is to capture the typical vessel boundary rather than extreme signal events. This trade-off may also explain why linearity behaved differently before and after contrast. Pre-contrast the vessel signal is weakest, and edge suppression by the RBC orientation artefact leaves little margin above the background, so the higher contrast of the MIP helps keep the profile distinguishable across vessels of different size. Once the lumen is filled with contrast agent, signal is no longer limiting and the outlier sensitivity of a maximum statistic becomes the dominant source of error, so the averaging of the MeIP yields the tighter relationship with the ground truth. Consistent with this, MeIP produced higher $R^2$ than MIP for all three model-based metrics post-contrast, whereas at peak contrast the two projections remained comparable, with MIP higher for two of the three.

The difference between MIP and MeIP was especially apparent in the binary mask based results (Table~\ref{tab:accuracy_linearity}). For MeIP, the slope declined more gradually and remained at 0.86 even at BW-5 post-contrast. Whereas for MIP, the slope decreased steeply with increasing threshold, dropping from 0.92 at BW-1 to 0.54 at BW-5 post-contrast. These findings suggest that MeIP produces a more uniform intensity distribution across the vessel profile, making the binarisation result less sensitive to the exact threshold value. The CoV data (Table~\ref{tab:accuracy_linearity}) reinforced this as the MIP-based binary mask CoV increased substantially with threshold, whereas the MeIP-based CoV remained relatively stable. This behaviour follows from how each projection responds to variance in the angiogram signal. The MIP retains only the brightest pixel along each A-scan, so it is a maximum statistic and follows the largest excursion in the signal, and the projected profile therefore varies from one time sample to the next. A fixed threshold applied to a varying profile returns a varying width, and the higher the threshold the nearer the peak that width is measured, where a given change in intensity displaces the threshold crossing further. The CoV consequently grows with threshold. MeIP instead averages along each A-scan, which suppresses this variance and produces a broader, smoother profile, leaving the CoV largely unchanged as the threshold increases. For model-based methods, MeIP also showed better performance. The $w_{1/e^2,\,\mathrm{G}}$ using MIP slightly overestimated vessel diameter post-contrast with $m = 1.12$, while MeIP reduced this to $m = 1.04$, a much closer agreement with the ground truth. These results suggest that while MIP may be preferred for visualisation purposes where contrast and SNR are important~\cite{hormelMaximumValueProjection2018}, MeIP should be the preferred projection method for quantitative diameter measurements.

\subsection{Influence of angiogram intensity on diameter measurements}
\label{subsec:influence_of_intensity}
The functional flicker stimulation experiment provided a direct demonstration of how changes in angiogram signal intensity can affect diameter measurements based on the method used. During the DyC-OCT acquisition, the angiogram intensity declined significantly across all five vessels (Fig.~\ref{fig:flicker}g,h), with decreases ranging from $-13.6\%$ to $-19.8\%$ (Table~\ref{tab:flicker}). The cause of this decline was not established with certainty, although inspection of the volumetric data ruled out the most important confounds. There was no out-of-plane motion or breathing artefact, and the image quality was consistent throughout the scan. We suspect a combination of rapid corneal drying and/or cataract formation over the course of the scan could have caused the decline in intensity. More importantly, the cause of the decline is not essential to our argument for this study, what matters is that the intensity decline biases the binary mask based diameter measurement while the model-based diameter remains unaffected. 

The model-based methods correctly detected the expected vasodilation, with the Gaussian model showing significant diameter increases in four of five vessels ($+15.5\%$ overall in the arteries against $+3.1\%$ in the veins) and the GG model showing significant increases in four of five vessels ($+14.0\%$ and $+1.4\%$ respectively; Table~\ref{tab:flicker}). The binary mask based method, however, reported significant diameter decreases in four of five vessels, suggesting vasoconstriction which is a paradoxical result. The reason behind this behaviour could be understood by observing how each method responds to declining signal intensity (Fig.~\ref{fig:scatter_intensity}). For the binary mask based method, the threshold is fixed and diameter is determined by the number of pixels that exceed this threshold. When the overall angiogram intensity drops, fewer pixels at the vessel edges exceed the threshold, and the vessel appears to shrink even if the physical vessel diameter has increased. This is confirmed by the strong positive correlation between measured diameter and angiogram intensity for the binary mask based method ($\rho = 0.526$, $p < 0.0005$). In contrast, the model-based methods showed weak overall correlations with intensity ($\rho = -0.101$ for Gaussian, $\rho = -0.098$ for GG), indicating that the fitted diameter was effectively independent of the signal intensity. Among the methods, the GG model also produced the lowest standard deviations across the five vessels in the flicker experiment (Fig.~\ref{fig:flicker}d--f), consistent with its low CoV post-contrast (Table~\ref{tab:accuracy_linearity}), indicating that it is the most repeatable of the metrics tested. Together these results demonstrated that the model-based diameter calculation provided more accurate and repeatable measurements, and recovered physiologically meaningful vascular responses that binary mask based methods misrepresented.

These findings have important implications for longitudinal studies and any experiment where angiogram intensity may vary between or within imaging sessions. In longitudinal mouse eye imaging, the angiogram signal can vary between and within sessions as anaesthesia alters retinal perfusion and ocular physiology~\cite{wormRetinalOCTParameters2026,fengHighresolutionStructuralFunctional2023}, with additional factors such as cataract formation, and corneal drying if tear drops are not applied, further reducing OCT intensity over time~\cite{steinEffectCornealDrying2006,mclellanOpticalCoherenceTomography2012}. If a binary mask based method is used under such conditions, observed changes in vessel diameter may reflect intensity fluctuations rather than true vascular changes. In contrast, model-based methods may be less sensitive to these intensity variations, as they characterise the shape of the intensity profile rather than relying primarily on absolute intensity values to define vessel boundaries.

\subsection{Threshold selection}
\label{subsec:threshold_selection}
The results presented here used a basic, global threshold for the binarisation of the angiogram data. There are numerous other approaches for binarising the angiogram, which will most likely each behave differently. For example, angiogram signals can be enhanced with vesselness filters~\cite{frangiMultiscaleVesselEnhancement1998} and dynamic thresholds can be used to binarise images with varying SNR. When applying a vesselness filter, it is possible to enhance the shape of the angiogram and omit noise, but it is also possible for such enhancement to overcompensate the signal, leading to larger than expected vessel appearance. Measurements obtained with different enhancement filters have been shown to differ substantially and are not interchangeable~\cite{hongEffectVesselEnhancement2020}. Dynamic thresholding can also yield a more uniform angiogram appearance, but similarly risks over- or under-compensating vessels. The threshold chosen is known to alter OCTA quantitative metrics systematically, with automatic thresholds tending to underestimate relative to a fixed threshold~\cite{mehtaImpactBinarizationThresholding2019,arrigoImpactDifferentThresholds2021}. Further analysis of these other methods would be required to determine their efficacy in returning accurate and precise diameter values.

The key observation from the present work is that any thresholding approach, regardless of how the threshold is determined, shares the same fundamental limitation that the measured diameter depends on the absolute signal intensity at the vessel edges. As demonstrated by the flicker stimulation experiment, this intensity dependence can introduce artefactual diameter changes when the angiogram signal varies over time or between sessions. A more adaptive thresholding strategy may reduce variability across different regions of a single image, but it would not eliminate the sensitivity to temporal intensity fluctuations that we observed. The model-based approach avoids this limitation entirely by deriving diameter from the shape of the intensity profile rather than from an intensity cut-off, making it a fundamentally different and more robust alternative for quantitative diameter analysis.

\subsection{Best choices under different conditions}
\label{subsec:best_choices}
Under the conditions where a contrast agent is available and accuracy is the primary goal, the $w_{1/e^2,\,\mathrm{GG}}$ metric using MeIP provides the best results post-contrast. This combination achieved a slope of $m = 1.001$ with $R^2 = 0.97$ and a CoV of 0.09, yielding the best accuracy, linearity, and precision of all tested methods. If only pre-contrast data is available, the $w_{1/e^2,\,\mathrm{GG}}$ metric still provides the best linearity among all tested methods ($R^2 = 0.81$ for MIP, $R^2 = 0.76$ for MeIP), though the MIP and MeIP slope of $m = 0.64$ and $m = 0.71$ indicates that diameter could be underestimated by approximately 35\% and 30\% respectively. In this case, a correction factor could be applied if absolute diameter values are needed, though relative comparisons between vessels or timepoints would remain valid without correction.

When only MIP data is available, which is the case for many commercial OCTA machines and existing datasets~\cite{sampsonTowardsStandardizingRetinal2022,untrachtTowardsStandardisingRetinal2024}, the $w_{1/e^2,\,\mathrm{GG}}$ still provides good post-contrast accuracy ($m = 1.02$, $R^2 = 0.95$). For binary mask based analysis of MIP data, only the lowest threshold (BW-1) post-contrast achieved a slope near unity ($m = 0.92$), and the slope degraded rapidly with increasing threshold values. This highlights that if binarisation must be used with MIP data, a lower threshold value is essential, though this may come at the cost of increased noise sensitivity.

The improved linearity and repeatability of the model-based approaches come at the cost of computation time, as fitting a model is considerably more expensive than binarisation. In our testing, fitting the generalised Gaussian was approximately 40\% slower than the standard Gaussian. This difference is hardware dependent, and because the fits at each cross-section are independent, the computation could be parallelised and accelerated with appropriate hardware. If computation time is a limiting factor and model fitting is not practical, MeIP with a moderate threshold (e.g.\ BW-3) provides a reasonable alternative. At post-contrast, BW-3 using MeIP achieved $m = 0.96$ and $R^2 = 0.81$, which is acceptable for many applications. MeIP-based binary mask based analysis showed more stable CoV values across all thresholds compared to MIP, making the results less sensitive to the exact threshold choice. This stability makes MeIP the preferred projection for binary mask based OCTA.

\subsection{Literature comparison}
\label{subsec:literature_comparison}

Hormel et~al.~\cite{hormelMaximumValueProjection2018} showed that MIP produces en face OCTA images with higher SNR and contrast compared to MeIP, and concluded that MIP should be preferred for OCTA visualisation. Our results reveal a different picture when it comes to quantitative vessel diameter measurement. We found that MeIP consistently yielded more accurate and precise diameter values than MIP across both model-based and binary mask based approaches, with the exception of pre-contrast linearity, where MIP performed better. This suggests that the projection method best suited for visualisation is not necessarily the best suited for quantification, and that the choice of projection should depend on the intended analysis.

The challenge of standardising OCTA image analysis has been highlighted by Sampson et~al.~\cite{sampsonTowardsStandardizingRetinal2022} and Untracht et~al.~\cite{untrachtTowardsStandardisingRetinal2024}, both of whom emphasised the need for transparent, reproducible analysis pipelines. Our findings add to this discussion by showing that the choice of diameter measurement method introduces a source of variability that is at least as significant as the choice of binarisation threshold. The model-based approach presented here offers a path toward more standardised vessel diameter quantification, as it reduces the dependence on selected thresholds and angiogram intensity, and provides measurements that are more robust across imaging conditions.

The RBC orientation artefact and its effect on OCTA vessel appearance has been described by Cimalla et al.~\cite{cimallaShearFlowinducedOptical2011} and Bernucci et~al.~\cite{bernucciInvestigationArtifactsRetinal2018}, who used Intralipid contrast to demonstrate that the hourglass pattern in vessel cross-sections is caused by the directional scattering properties of RBCs rather than by true variations in flow. Our work extends this observation by quantifying the degree to which this artefact affects diameter measurements and by proposing model-based fitting as a practical correction strategy. While contrast enhancement with Intralipid is effective for revealing the true vessel diameter, it requires a contrast injection and is not yet applicable in clinical settings. Using a calibrated model-based approach provides a way to partially recover the underestimated diameter from native OCTA data without the need for a contrast agent.

Recent work by Son et~al.~\cite{sonFunctionalOCTReveals2024} used functional OCT to image flicker-evoked vasodilation in the mouse retina and reported that vasodilation is anisotropic, with larger percent changes in the axial direction ($\sim18\%$) than the lateral direction ($\sim11\%$). The authors did not describe how vessel diameters were quantified from their B-scans. It is worth considering whether the lateral diameter measurements may have been affected by the intensity-dependent artefacts observed in the present work. In the lateral direction, the RBC orientation artefact suppresses signal at vessel edges, and depending on how their vessel boundary was defined, this artefact could lead to an underestimation of the lateral diameter. If this were the case, the apparent difference between axial and lateral vasodilation could be influenced by the measurement methodology in addition to the biomechanical anisotropy discussed by the authors. It would be interesting to apply the model-based approach presented here to such data to determine whether their reported anisotropy changes when the lateral diameter is quantified differently.

The magnitude of vasodilation detected by the generalised Gaussian model ranged from $+3.1\%$ to $+28.4\%$ across the vessels with a significant response (Table~\ref{tab:flicker}). The arteries (Vessels 1, 2, and 4) showed increases of $+3.5\%$, $+10.1\%$, and $+28.4\%$, while the veins (Vessels 3 and 5) showed a smaller response, with $+3.1\%$ in Vessel 5 and no significant change in Vessel 3. The standard Gaussian gave a broadly similar picture, with arterial increases of $+8.2\%$, $+9.9\%$, and $+28.5\%$ and a venous increase of $+4.9\%$ in Vessel 5. Table~\ref{tab:flicker_comparison} summarises the stimulation protocols and vascular responses reported by previous flicker stimulation studies alongside our own. The listed human studies span several measurement techniques, from fundus photography~\cite{formazDiffuseLuminanceFlicker1997} through retinal vessel analysers~\cite{polakInfluenceFlickerFrequency2002,garhoferDiffuseLuminanceFlicker2004,nagelFlickerObservationLight2004} to Doppler OCT~\cite{aschingerEffectofDiffuse2017} and OCTA~\cite{kallabPlexusSpecificEffect2020}. Despite these methodological differences, the reported human arterial responses cluster between $+3\%$ and $+8\%$, and venous responses between $+2\%$ and $+10\%$. Notably, these studies report arterial and venous responses of broadly comparable magnitude, and in some cases the venous response is the larger of the two, for example $+9.8\%$ venous versus $+8.0\%$ arterial in Aschinger et~al.~\cite{aschingerEffectofDiffuse2017}. This contrasts with our measurements, in which the arterial response was on average larger than the venous response. Whether this asymmetry reflects a genuine feature of the murine flicker response, the small number of vessels sampled, or a greater sensitivity of the model-based fit to the smooth-muscle-driven dilation of arteries remains to be established. Polak et~al.~\cite{polakInfluenceFlickerFrequency2002} additionally showed that the diameter response is relatively stable across flicker frequencies from 4 to 40~Hz, suggesting that frequency differences between studies are unlikely to explain the variation.

The magnitudes we observed are generally larger than those reported in the human studies, though direct comparison is difficult given the differences in experimental conditions. Species, anaesthesia, stimulation frequency and wavelength, stimulus power and delivery method, stimulation duration, dark adaptation period, and imaging modality all vary across the reported studies, and each of these factors could influence the measured vascular response. Among the studies listed, Son et~al.~\cite{sonFunctionalOCTReveals2024} used the most similar setup to ours. They used the same mouse strain (C57BL/6J), the same anaesthetic (ketamine/xylazine), a comparable flicker stimulus (10~Hz, 504~nm), however a different stimulation period of 5~s. Because the diameters reported here are lateral, their lateral values offer the closest comparison: they reported $+10.8\%$ in the arteries and $+8.9\%$ in the veins, against $+14.0\%$ and $+1.4\%$ respectively for the GG model in this study. The arterial responses are therefore of comparable magnitude, whereas our venous response is substantially smaller. A systematic comparison of these parameters could help determine which factors are most influential to the results observed in both studies. Importantly, regardless of the absolute magnitude, the fact that we were able to detect vasodilation at all depended on the measurement method, because binary mask based analysis not only failed to detect vasodilation but reported vasoconstriction which would have led to fundamentally incorrect physiological conclusions.

Previous studies in mice have measured flicker-evoked retinal vessel diameter changes mostly using dynamic vessel analysis (DVA), which is a fundus camera-based technique adapted from human retinal vessel analysers. Albanna et~al.~\cite{albannaNonInvasiveEvaluation2018} demonstrated the feasibility of DVA in the mouse eye using 12.5~Hz flicker at 530~nm but found that quantitative arterial recordings were possible only in a minority of animals, while venous dilation was mostly below 1.5\% post-stimulation. Another study by Neumaier et~al.~\cite{neumaierEvaluationFlickerInduced2021} used the same setup in wildtype and Cav2.3-deficient mice with overnight dark adaptation and reported median arterial and venous diameter changes of only $+1.3\%$ and $+1.2\%$, respectively, in the wildtype group. These values are substantially smaller than those we observed, which could be due to the limited spatial resolution and signal-to-noise ratio of DVA when applied to the smaller mouse eye rather than being a genuine physiological response. The much larger responses we observed may therefore partly result from the higher lateral resolution of OCT combined with the model-based fitting approach, which captures vessel boundary shifts that could fall below the detection limit of DVA. In addition to these diameter-based studies, Rai et~al.~\cite{raiEffectsFlickeringLight2025} recently reported that 12~Hz flicker stimulation significantly increased both arterial and venous retinal blood flow in light-adapted C57BL/6J mice. An earlier OCT-based study by Radhakrishnan and Srinivasan~\cite{radhakrishnanMultiparametricOpticalCoherence2013} measured arterial dilation in the rat inner retina during a 10~s, 12~Hz light stimulus, reporting a 95\% confidence interval of $+3.3\%$ to $+11.0\%$ ($p < 0.005$) for diameters quantified manually along the arterial tree. Our arterial response ($+14.0\%$ for the GG model) sits just above the upper limit of that interval, which is a reasonable correspondence given the difference in species. 

\begin{table}[htbp]
\centering
\caption{Comparison of flicker stimulation protocols and reported diameter changes across studies. H = human, M = mouse, R = rat. DA = dark adaptation. Lat.\ = lateral, ax.\ = axial.}
\label{tab:flicker_comparison}
\resizebox{\textwidth}{!}{%
\begin{tabular}{lccccccc}
\hline
 & \textbf{Stimulation} & \textbf{Stimulation} & \textbf{Stimulation} & \textbf{DA} & \multicolumn{2}{c}{\textbf{Diameter Change (\%)}} \\
\cline{6-7}
\textbf{Study} & \textbf{Freq.~(Hz)} & \textbf{Light Source} & \textbf{Duration~(s)} & \textbf{(h)} & \textbf{Artery} & \textbf{Vein} \\
\hline
Formaz~\cite{formazDiffuseLuminanceFlicker1997} (H) & $10$ & Xenon arc, $5.4~log\,Td^{\S}$ & $60$ & --- & $+4.2$ & $+2.7$ \\
Polak~\cite{polakInfluenceFlickerFrequency2002} (H) & $8$ & $<550~nm$, $\sim150\,\mu\mathrm{W},cm^{-2}$$^{\dagger}$ & $60$ & --- & $+3.5$ & $+1.9$ \\
Garhofer~\cite{garhoferDiffuseLuminanceFlicker2004} (H) & $8$ & $<550~nm$, $300\,\mu\mathrm{W},cm^{-2}$$^{\dagger}$ & $60$ & --- & $+3.8$ & $+2.8$ \\
Nagel~\cite{nagelFlickerObservationLight2004} (H) & $12.5$ & $530-600~nm$, $\sim196\,\mu\mathrm{W},cm^{-2}$$^{\dagger}$ & $20$ & --- & $+6.9$ & $+6.5$ \\
Aschinger~\cite{aschingerEffectofDiffuse2017} (H) & $12.5$ & $530-600~nm$, $\sim196\,\mu\mathrm{W},cm^{-2}$$^{\dagger}$ & $120$ & --- & $+8.0$ & $+9.8$ \\
Kallab~\cite{kallabPlexusSpecificEffect2020} (H) & $7$ & White, $\sim450~lux$ & $^{*}$ & --- & $+3.7$ & $+4.5$ \\
Albanna~\cite{albannaNonInvasiveEvaluation2018} (M) & $12.5$ & $530~nm$, $30~lux$ & $20$ & $12$ & ---$^{\#}$ & $<1.5^{\#}$ \\
Neumaier~\cite{neumaierEvaluationFlickerInduced2021} (M) & $12.5$ & $530~nm$, $30~lux$ & $20$ & $12$ & $+1.3$ & $+1.2$ \\
Radhakrishnan~\cite{radhakrishnanMultiparametricOpticalCoherence2013} (R) & $12$ & Hg:Xe, $2\,\mathrm{mW},cm^{-2}$$^{\dagger}$ & $10$ & --- & $+3.3$ to $+11.0^{\parallel}$ & ---$^{\parallel}$ \\
Son~\cite{sonFunctionalOCTReveals2024} (M) & $10$ & $504~nm$, $\sim460~lux$ & $5$ & $1-2$ & lat.\ $+10.8$, ax.\ $+18.2$ & lat.\ $+8.9$, ax.\ $+14.3$ \\
Present study (M) & $10$ & $502~nm$, $26.6\,\mu\mathrm{W}$$^{\ddagger}$ & Cont. & $2$ & $+3.5$ to $+28.4^{\P}$ & $+3.1^{\P}$ \\
\hline
\multicolumn{7}{l}{\footnotesize $^{\S}$Diameter measured from fundus photographs taken post-stimulation; retinal luminance in trolands.} \\
\multicolumn{7}{l}{\footnotesize $^{\#}$DVA feasibility study; arterial recordings were not reliably obtainable in the mouse eye.} \\
\multicolumn{7}{l}{\footnotesize $^{\P}$Values from individual vessels with a statistically significant response; other studies report group means.} \\
\multicolumn{7}{l}{\footnotesize $^{\parallel}$95\% confidence interval, not a range of per-vessel means. Diameter was quantified along the arterial tree only, so no venous value is reported.} \\
\multicolumn{7}{l}{\footnotesize $^{*}$Flicker applied during OCTA scan acquisition. $^{\dagger}$Retinal irradiance. $^{\ddagger}$Power at cornea.} \\
\end{tabular}%
}
\end{table}

\subsection{Future direction}
\label{subsec:future_direction}
A future direction of this work is to move toward fully automated segmentation and calculation of vessel diameters from en face angiogram maps. While ROIs were manually selected in this work, it is also possible to automatically segment angiographic data using traditional skeletonisation methods, which are commonly used for automated OCTA analysis. From a skeletonised image, vessel angle can be computed for further extraction of vascular cross-sections perpendicular to the vessel axis. These cross-sections could then be processed using the model-based methods described here to enable fully automated analysis of 3D OCTA data sets, rather than the semi-automated approach shown here for 2D OCTA data. This type of model-based approach is particularly important for comparing 3D volumes acquired at different time points in the mouse eye, as there are often intensity losses caused by drying of the eye, cataract formation, etc. As we have shown in this work, even low levels of intensity loss occurring over a short time window can significantly affect the measured vessel diameter when using a binary mask based approach, but not when using the model based approach presented here.

Additionally, this work focused on the larger vessels of the superficial vascular plexus. Extending the model-based approach to the intermediate and deep capillary plexuses, where vessels are smaller and signal levels are lower, would be valuable for studying microvascular changes in diseases. For very small vessels approaching the lateral resolution limit of the imaging system, the measured intensity profile is a convolution of the true vessel profile with the point spread function (PSF) of the incident beam on the retina. Since the lateral resolution of OCT is typically worse than the axial resolution and is determined by the beam optics rather than the source bandwidth, this convolution becomes increasingly significant for smaller vessels. If the PSF were known, a deconvolution approach could be used to recover the true vessel profile before model fitting, potentially improving accuracy for capillary-scale vessels. However, characterising the retinal PSF in vivo remains challenging, as it depends on the ocular optics of each individual eye and may vary with retinal depth and eccentricity.
